% BHU PYQ -- Manually transcribed & error-corrected
\documentclass[11pt,a4paper]{article}

\usepackage[a4paper,top=2.5cm,bottom=2.5cm,left=2.8cm,right=2.8cm,headheight=14pt]{geometry}
\usepackage{amsmath,amssymb}
\usepackage[T1]{fontenc}
\usepackage[utf8]{inputenc}
\usepackage{lmodern}
\usepackage{microtype}
\usepackage{fancyhdr}
\usepackage{enumitem}
\usepackage{hyperref}
\usepackage{lastpage}
\usepackage{parskip}

\hypersetup{colorlinks=true,urlcolor=black,linkcolor=black,
  pdftitle={PHYSICS: Ancillary Physics-II -- BHU PYQ},pdfauthor={Banaras Hindu University}}

\pagestyle{fancy}\fancyhf{}
\renewcommand{\headrulewidth}{0.4pt}\renewcommand{\footrulewidth}{0.4pt}
\fancyhead[L]{\small   \textit{Banaras Hindu University}}
\fancyhead[R]{\small   \textit{Ancillary Physics-II}}
\fancyfoot[C]{\small Page  \thepage\ of \pageref{LastPage}}


\newlist{parts}{enumerate}{2}
\setlist[parts,1]{label=(\alph*),leftmargin=*,itemsep=5pt,topsep=3pt}
\setlist[parts,2]{label=(\roman*),leftmargin=*,itemsep=3pt,topsep=2pt}

\newcommand{\pts}[1]{\hfill   \textit{[#1]}}


\begin{document}


\begin{center}
  {\large\bfseries BANARAS HINDU UNIVERSITY}\\[2pt]
  {
\normalsize B.Sc. (Hons.) Semester IV Examination, 2013-14}\\[6pt]
  \rule{0.8\linewidth}{0.5pt}\\[6pt]
  {\large\bfseries Paper: BSC07A / BP Anc.-II: Ancillary Physics-II}\\[4pt]
  {
\normalsize Time: 3 Hours\hfill Maximum Marks: 70}
\end{center}
\rule{\linewidth}{0.4pt}
\smallskip



\noindent   \textit{Instructions: Answer   \textbf{five} questions including Question~1 which is compulsory. Symbols have their usual meanings.}
\bigskip




\noindent   \textbf{Question 1.} Answer any   \textbf{six} of the following: \pts{6 $ \times$ 3 = 18}

\begin{parts}
  \item What do you mean by intensive and extensive properties?
  \item Explain the dual nature of light with examples.
  \item How can diffusion and Brownian motion be explained on the basis of random walk?
  \item What will be the value of entropy at $T = 0\, \text{K}$? Will a Carnot cycle work at this temperature?
  \item Green light of wavelength $5100\, \text{\AA}$ from a narrow slit is incident on a double slit. If the overall separation of 10 fringes on a screen $200\, \text{cm}$ away is $2\, \text{cm}$, find the slit separation.
  \item Explain various uses of polarization of light.
  \item Define the power of a lens. Mention the formula for the power of a combination of two thin lenses.
  \item What do you mean by thermal equilibrium? Which law of thermodynamics defines it?
\end{parts}

\medskip\rule{\linewidth}{0.2pt}




\noindent   \textbf{Question 2.}

\begin{parts}
  \item Explain various types of defects in the human eye. How can these defects be removed using different lenses? \pts{5}
  \item State and explain Fermat's principle of extremum path with an example. \pts{3}
  \item Explain the differences between the Fresnel and Fraunhofer classes of diffraction. \pts{5}
\end{parts}

\medskip\rule{\linewidth}{0.2pt}




\noindent   \textbf{Question 3.}

\begin{parts}
  \item Define the resolving power of an optical instrument. Obtain an expression for the resolving power of a telescope. \pts{7}
  \item A point source of light is located $20\, \text{cm}$ in front of a convex mirror with a focal length of $15\, \text{cm}$. Determine the character and position of the image point. \pts{6}
\end{parts}

\medskip\rule{\linewidth}{0.2pt}




\noindent   \textbf{Question 4.}

\begin{parts}
  \item Calculate the focal length of a plano-convex lens of which the radius of curvature of the curved surface is $40\, \text{cm}$ and the refractive index is $\mu = 1.5$. \pts{5}
  \item Explain the construction and working of a compound microscope using a proper ray diagram. \pts{8}
\end{parts}

\medskip\rule{\linewidth}{0.2pt}




\noindent   \textbf{Question 5.}

\begin{parts}
  \item Establish the relations for thermodynamic variables $S$, $T$, $P$, and $V$ in terms of thermodynamic potentials $U$, $F$, $H$, and $G$, where the symbols have their usual meanings. \pts{8}
  \item What do you mean by an isothermal process? Show that in an isothermal expansion of an ideal gas, $\Delta U = 0$ and $\Delta H = 0$. \pts{5}
\end{parts}

\medskip\rule{\linewidth}{0.2pt}




\noindent   \textbf{Question 6.}

\begin{parts}
  \item Discuss the working of a Carnot engine and sketch its cycle on a $P$-$V$ diagram. Explain the four processes involved. \pts{6}
  \item The efficiency of a Carnot engine can be increased in two ways: (i) by decreasing the temperature of the sink $T_2$ while keeping the temperature of the source $T_1$ constant, or (ii) by increasing the temperature of the source $T_1$ while keeping the temperature of the sink $T_2$ constant. Which of these methods is more effective and why? \pts{7}
\end{parts}

\medskip\rule{\linewidth}{0.2pt}




\noindent   \textbf{Question 7.}

\begin{parts}
  \item What is the difference between effusion and diffusion? \pts{3}
  \item State Graham's law of diffusion. \pts{3}
  \item What do you mean by drift velocity and mean free path of a gas molecule? \pts{4}
  \item Why are water droplets formed when a bicycle tube bursts? Explain. \pts{3}
\end{parts}

\vfill
\rule{\linewidth}{0.4pt}

\begin{center}\small
  \href{https://bhu-exam.vercel.app/}{Click here to visit our website}\\[4pt]
  \href{https://me-aryan.vercel.app/}{Click here to visit our contributor}\\[4pt]
  \href{https://quashan-me.vercel.app/}{Click here to visit our contributor}
\end{center}

\end{document}
